
\subsubsection{6.1 Interpretation of Findings}

\textbf{Framework validation.} The most secure finding is that the structural efficiency measurement framework executes correctly end-to-end. All five modules produce interpretable outputs on real DeepSeek-Coder-7B-instruct-v1.5 checkpoints without errors. This validates the measurement infrastructure as a tool that can be applied to checkpoint archives from any GRPO or DPO training run.

\textbf{Observation of near-equal SEP.} The h-m1 analysis produced SEP ≈ 0.237 for both GRPO and DPO. This result, taken at face value, does not support the hypothesis that execution reward selectively concentrates policy movement on control-flow and data-flow AST nodes. However, three competing interpretations apply:

(a) \textit{The mechanism claim may be false even if the existence claim holds.} GRPO may produce more total structural change in absolute terms — consistent with the h-e1 raw edit distance observation on synthetic data — without proportionally concentrating edits on semantic nodes. Under this interpretation, execution reward produces broadly more aggressive restructuring rather than selective semantic node targeting.

(b) \textit{Checkpoint aliasing may mask a real SEP difference.} With n_eff ≈ 2, the Mann-Whitney test has essentially no statistical power. A corrected run with 10 or more diverse checkpoints might reveal SEP_GRPO > SEP_DPO. This interpretation cannot be excluded on the basis of the current data.

(c) \textit{Scale-level SEP equilibration.} At 7B parameter scale with the specific KL budget used, both methods may converge to similar proportional distributions over node types regardless of absolute edit volume. This would be a scale-specific finding.

Distinguishing these interpretations requires a corrected experimental run as described in Section 6.3.

\textbf{Checkpoint aliasing as a methodological finding.} The identification of checkpoint aliasing as a failure mode that silently corrupts nominally large analyses has value independent of any finding about GRPO versus DPO. The confound is invisible from benchmark metrics and requires structural analysis to detect. Any researcher conducting checkpoint-comparison analysis in RL fine-tuning — for code, mathematics, or general NLP tasks — should add checkpoint diversity pre-flight checks to their pipeline.

\subsubsection{6.2 Limitations}

\textbf{L1: Synthetic data in h-e1 gate evaluation.} The proof-of-concept experiment used hand-crafted code completions with controlled structural properties, not real GRPO training outputs. The bootstrap CI [4.65, 8.73] and mean differential 6.50 are valid measurements of the infrastructure's detection capability given those inputs; they cannot be cited as empirical evidence that real GRPO training produces higher structural efficiency than DPO.

\textbf{L2: Underpowered SEP analysis.} The h-m1 SEP comparison is entirely underpowered (n_eff ≈ 2). The Mann-Whitney test result (p = 0.4248) and the near-zero effect size (−0.0072) cannot be interpreted as evidence for or against the mechanism hypothesis.

\textbf{L3: Single base model and language.} All experiments use DeepSeek-Coder-7B-instruct-v1.5 on Python. The structural efficiency framework is designed to be model-agnostic and applicable to any AST-parseable language, but empirical generalization across scales and languages has not been tested.

\textbf{L4: DPO preference pairs not execution-oracle labeled.} The h-e1 DPO training used hard-coded \texttt{return None} as the rejected completion rather than genuine model-generated failed solutions labeled by evalplus. The DPO training condition in the proof-of-concept does not represent execution-oracle DPO as specified. This may suppress DPO's structural activity and inflate the apparent raw edit distance difference in the synthetic results.

\textbf{L5: KL tolerance inflation.} The proof-of-concept run used KL tolerance 0.15 rather than the specified ±5\% (0.05). KL-matched pairs are therefore less precisely controlled than intended.

\textbf{L6: SEP not validated against functional outcomes.} The structural efficiency metric and SEP have not been validated against functional correctness measures such as pass@1, expected calibration error, or out-of-distribution transfer. The framework measures structural activity of policy movement, not the correctness of that movement.

\subsubsection{6.3 Corrected Experimental Protocol}

For a valid empirical test of the mechanism hypothesis:

\begin{enumerate}
\item \textbf{Enforce checkpoint diversity:} Use \texttt{save_steps=100} for a 1000-step run, yielding 10 GRPO and 10 DPO checkpoints. Assert \texttt{len(unique_checkpoints) >= 10} before proceeding to h-m1 analysis.
\item \textbf{Replace synthetic data:} Replace hand-coded completions in \texttt{smoke_test_experiment.py} with real GRPOTrainer outputs evaluated by evalplus.
\item \textbf{Fix DPO pairs:} Implement execution-oracle DPO pair generation: sample model outputs, label by evalplus execution (passing = preferred, failing = rejected).
\item \textbf{Use specified KL tolerance:} Use tolerance = 0.05 (±5\%). With 10+ real checkpoints, sufficient pairs will be found.
\item \textbf{Reduce scope:} Confirm the h-e1 → h-m1 chain produces valid results before attempting downstream sub-hypotheses (h-m2 through h-m4).
\end{enumerate}

\subsubsection{6.4 Broader Impact}

The structural efficiency framework provides a diagnostic for comparing post-training methods at the level of policy movement structure rather than benchmark outcomes. This may reduce reliance on pass@1 as the sole comparison criterion and enable more targeted analysis of what different alignment methods learn. The checkpoint aliasing finding provides a concrete, low-cost safeguard for any team conducting checkpoint-comparison analysis of RL fine-tuning.

The structural efficiency metric should not be used as a training signal until its relationship to functional correctness and out-of-distribution generalization is empirically established, as optimization pressure toward high SEP without functional correctness constraints could produce structurally complex but incorrect code.
